\chapter{Liposuction}
\index{Liposuction}

\section*{Definition and Overview}
Liposuction, also known as lipoplasty or suction-assisted lipectomy, is a surgical procedure that uses a suction technique to remove fat from specific areas of the body, such as the abdomen, hips, thighs, buttocks, arms, or neck. It is a cosmetic and reconstructive procedure designed to contour the body rather than serve as a weight-loss method or a treatment for obesity. Various techniques exist, including tumescent, ultrasound-assisted, and laser-assisted liposuction. The procedure aims to target localised fat deposits that are resistant to traditional diet and exercise regimes.

\section*{Epidemiology}
Liposuction is one of the most frequently performed cosmetic surgical procedures worldwide. In many countries and other Western nations, it consistently ranks in the top five cosmetic surgical procedures for both men and women. The demand has remained high over recent decades, with an increasing number of males seeking the procedure, particularly for gynaecomastia and abdominal contouring. However, patient safety guidelines have become more stringent in many countries, limiting the volume of fat removed in a single session to reduce systemic risk and improve clinical outcomes.

\section*{Aetiology and Pathophysiology}
Liposuction physically disrupts and removes subcutaneous adipose tissue, which is relatively resistant to diet and exercise due to localized hormonal and genetic factors.
\begin{itemize}
    \item \textbf{Subcutaneous fat deposition:} distribution is governed by genetics, age, and hormonal factors, forming localized pockets.
    \item \textbf{Tumescent fluid infiltration:} a solution of saline, adrenaline (vasoconstrictor), and lignocaine (local anaesthetic) is injected to minimise blood loss and provide postoperative analgesia.
    \item \textbf{Mechanical disruption:} a cannula inserted through small incisions is moved back and forth to loosen adipose tissue.
    \item \textbf{Suction and fluid shifts:} high-vacuum pressure extracts the fat, leading to fluid shifts and third-spacing that require careful monitoring.
\end{itemize}

\section*{Clinical Features}
Liposuction is elective, and patients undergo detailed preoperative assessment to determine anatomical suitability and realistic expectations.
\begin{itemize}
    \item Localised deposits of subcutaneous fat that are disproportionate to the rest of the body
    \item Firm skin with good elasticity, which is essential for the skin to contract smoothly over the newly contoured areas
    \item Minimal skin laxity or sagging (excess skin may require excisional surgery like abdominoplasty instead)
    \item Postoperative signs including transient swelling, bruising, numbness, and mild discomfort in the treated areas
\end{itemize}

\section*{Differential Diagnosis}
The clinical decision involves distinguishing patients suitable for liposuction from those needing alternative procedures or medical treatments.
\begin{itemize}
    \item \textbf{Obesity:} generalized excess adipose tissue, which is best managed with lifestyle modifications, pharmacotherapy, or bariatric surgery.
    \item \textbf{Skin laxity and ptosis:} excess, hanging skin (e.g., post-pregnancy or significant weight loss) that requires a lift or excision (e.g., tummy tuck) rather than liposuction.
    \item \textbf{Visceral fat accumulation:} intra-abdominal fat surrounding internal organs, which cannot be removed by liposuction.
    \item \textbf{Lipoedema:} a chronic adipose tissue disorder characterised by symmetrical fat deposition in the limbs, which may require specialised lymphatic-sparing liposuction.
    \item \textbf{Gynaecomastia vs.\ pseudogynaecomastia:} glandular breast tissue in males requires excision, whereas fatty tissue (pseudogynaecomastia) is suitable for liposuction.
\end{itemize}

\section*{When to Seek Medical Care}
Patients must be thoroughly educated on postoperative warning signs that indicate surgical or systemic complications.

\textbf{Contact your local emergency services for:}
\begin{itemize}
    \item Sudden shortness of breath, rapid breathing, or chest pain (suggesting pulmonary embolism or fat embolism)
    \item Severe dizziness, confusion, or loss of consciousness
    \item Rapidly spreading redness, extreme swelling, or high fever (suggesting severe infection or necrotising fasciitis)
    \item Profuse bleeding or rapidly expanding swelling (haematoma) at the incision sites
\end{itemize}

\section*{Diagnosis and Investigations}
Assessment is primarily clinical and structural.
\begin{itemize}
    \item \textbf{Preoperative clinical assessment:} evaluating skin elasticity, fat distribution, and medical history.
    \item \textbf{Routine preoperative blood tests:} full blood count, coagulation profile, and renal function to ensure surgical safety.
    \item \textbf{Cardiovascular screening:} ECG and other investigations for patients with cardiovascular risk factors or those undergoing general anaesthesia.
\end{itemize}

\section*{Management}
Postoperative care and adherence to guidelines are critical for optimal cosmetic and health outcomes.
\begin{itemize}
    \item \textbf{Compression garments:} wearing specialized elastic garments for several weeks to reduce swelling, support healing tissue, and assist skin contraction.
    \item \textbf{Pain management:} prescribed or over-the-counter oral analgesics to control postoperative pain and discomfort.
    \item \textbf{Fluid and electrolyte management:} monitoring fluid balance, especially during large-volume liposuction, to prevent volume depletion or overload.
    \item \textbf{Activity modification:} gradual return to light activities, avoiding strenuous exercise for several weeks as advised by the surgeon.
    \item \textbf{Lymphatic drainage massage:} gentle massage or physiotherapy to help resolve residual swelling and prevent contour irregularities.
\end{itemize}

\section*{Complications}
\begin{itemize}
    \item Contour irregularities, asymmetry, or dimpling of the skin
    \item Fluid accumulation (seromas) or haematoma formation at the surgical site
    \item Infection, ranging from localized cellulitis to life-threatening necrotising infections
    \item Lidocaine toxicity or fluid overload from the tumescent solution
    \item Fat embolism syndrome or deep vein thrombosis (DVT) and pulmonary embolism (PE)
    \item Skin necrosis or thermal burns from ultrasound/laser-assisted devices
\end{itemize}

\section*{Prognosis and Outcomes}
The long-term outcomes of liposuction are generally highly satisfactory, provided the patient maintains a stable weight. The removed fat cells do not regenerate, but remaining fat cells can enlarge if the patient gains weight. Swelling typically resolves over three to six months, revealing the final contour. Complication rates are low when the procedure is performed by a qualified specialist surgeon in an accredited facility.

\section*{Prevention and Self-Care}
\begin{itemize}
    \item Maintain a stable weight through a balanced diet and regular exercise to preserve cosmetic results.
    \item Wear the compression garment strictly as prescribed by the surgical team.
    \item Avoid smoking and alcohol for several weeks before and after surgery to promote optimal wound healing.
    \item Mobilise early by walking gently to prevent DVT and improve circulation.
    \item Keep the incision sites clean and dry, inspecting them daily for signs of infection.
\end{itemize}

\section*{Sources and Further Reading}
The following reputable organisations provide further information on liposuction, cosmetic surgery guidelines, and safety.
\begin{itemize}
    \item Australian Society of Plastic Surgeons. \textit{Liposuction procedure information and patient guide.} https://www.plasticsurgery.org.au/
    \item Mayo Clinic. \textit{Liposuction risks, results, and what you can expect.} https://www.mayoclinic.org/tests-procedures/liposuction/
    \item American Society of Plastic Surgeons. \textit{Liposuction cosmetic surgery details.} https://www.plasticsurgery.org/
    \item Healthdirect Australia. \textit{Cosmetic surgery: liposuction safety and considerations.} https://www.healthdirect.gov.au/cosmetic-surgery-liposuction
    \item Better Health Channel. \textit{Liposuction surgery overview.} https://www.betterhealth.vic.gov.au/health/conditionsandtreatments/liposuction
\end{itemize}
